% =============================================================================
%  Body of the worksheet "<<TITLE>>".
%  Only the content of the document belongs here -- no preamble, no
%  \begin{document}. Everything around is supplied by the header template.
%
%  \begin{task}[title=..., points=..., stars=..., bloom=...,
%               answer=..., lines=..., space=...]
%      answer = none | blank | lines | dots | box | grid   (default: none)
%      lines  = number of lines used by answer=lines and answer=dots
%      space  = height used by answer=blank, answer=box and answer=grid
%      stars  = 1 or 2 -- difficulty printed behind the number of the task
%      bloom  = note for the teacher, printed in the solution sheet only
%
%  Note: \end{code} and \end{lstlisting} have to start at the beginning of a
%  line, otherwise the indentation is printed as an extra line of the listing.
% =============================================================================

\begin{task}[title=Zahřívací otázka, points=2, bloom=porozumění,
             answer=lines, lines=3]
  Vysvětlete vlastními slovy, k čemu slouží cyklus \texttt{for}.

  \begin{solution}
    Slouží k opakovanému provedení bloku příkazů pro každý prvek sekvence.
  \end{solution}
\end{task}

\begin{task}[title=Rozbor kódu, points=3, bloom=analýza, answer=dots, lines=2]
  Co vypíše následující program?

  \begin{code}[Python]
for i in range(3):
    print(i * i)
\end{code}

  \begin{solution}
    Vypíše čísla 0, 1 a 4, každé na samostatný řádek.
  \end{solution}
\end{task}

\begin{task}[points=5, stars=1, bloom=tvorba, answer=box, space=5cm]
  Napište program, který vypíše všechna sudá čísla od 1 do 20.

  \begin{solution}
    \begin{code}[Python]
for i in range(2, 21, 2):
    print(i)
\end{code}
  \end{solution}
\end{task}
